\documentclass[journal, onecolumn]{IEEEtran}

% --- Codificacion y lenguaje -------------------------------------------------
\usepackage[utf8]{inputenc}
\usepackage[T1]{fontenc}
\usepackage[spanish, es-tabla]{babel}

% --- Margenes ----------------------------------------------------------------
\usepackage[top=2.5cm, bottom=2.5cm, left=2.5cm, right=2.5cm]{geometry}

\usepackage{listings}
\usepackage{xcolor}

\lstdefinestyle{ArduinoStyle}{
    language=C++,
    basicstyle=\ttfamily\small,
    keywordstyle=\color{blue},
    commentstyle=\color{green!60!black},
    stringstyle=\color{orange},
    numbers=left,
    numberstyle=\tiny,
    frame=single,
    breaklines=true,
    backgroundcolor=\color{gray!10}
}

% --- Graficos y figuras ------------------------------------------------------
\usepackage{graphicx}
\usepackage{caption}
\usepackage{float}

% --- Colores -----------------------------------------------------------------
\usepackage{xcolor}
\definecolor{azulTitulo}{RGB}{0,70,127}
\definecolor{codegreen}{rgb}{0,0.6,0}
\definecolor{codegray}{rgb}{0.5,0.5,0.5}
\definecolor{codepurple}{rgb}{0.58,0,0.82}
\definecolor{backcolour}{rgb}{0.95,0.95,0.92}

% --- Tablas ------------------------------------------------------------------
\usepackage{booktabs}
\usepackage{tabularx}
\usepackage{array}

% --- Hipervinculos -----------------------------------------------------------
\usepackage[hidelinks]{hyperref}

% --- Encabezado y pie de pagina ----------------------------------------------
\usepackage{fancyhdr}
\setlength{\headheight}{33pt}
\pagestyle{fancy}
\fancyhf{}
\lhead{%
  \begin{picture}(0,0)
    \put(5,10){\includegraphics[width=11mm]{pictures/escudo.png}}
  \end{picture}%
  \hspace{1mm}\textsc{\color{azulTitulo}Control Digital \hspace{90mm} Ing. Telecomunicaciones}%
}
\rhead{%
  \begin{picture}(125,0)
    \put(10,10){\includegraphics[width=35mm]{pictures/logo.png}}
  \end{picture}%
}
\cfoot{\thepage}
\renewcommand{\headrulewidth}{0.6pt}
\renewcommand{\sectionmark}[1]{}
\renewcommand{\subsectionmark}[1]{}

\fancypagestyle{plain}{
  \fancyhf{}
  \lhead{%
    \begin{picture}(0,0)
      \put(0,-5){\includegraphics[width=8.5mm]{pictures/escudo.png}}
    \end{picture}%
  }
  \rhead{%
    \begin{picture}(86,0)
      \put(0,-10){\includegraphics[width=30mm]{pictures/logo.png}}
    \end{picture}%
  }
  \cfoot{\thepage}
  \renewcommand{\headrulewidth}{0.6pt}
}

% --- Formato de secciones ----------------------------------------------------
\usepackage{titlesec}
\usepackage{parskip}

\titleformat{\section}
  {\color{azulTitulo}\large\bfseries}{\thesection.}{0.5em}{}
  [\color{azulTitulo}\titlerule]

\titleformat{\subsection}
  {\color{azulTitulo}\normalsize\bfseries}{\thesubsection}{0.5em}{}[]

\titleformat{\subsubsection}
  {\color{azulTitulo}\normalsize\itshape}{\thesubsubsection}{0.5em}{}[]

% --- Numeracion de secciones -------------------------------------------------
\renewcommand{\thesection}{\arabic{section}}
\renewcommand{\thesubsection}{\thesection.\arabic{subsection}}
\renewcommand{\thesubsubsection}{\thesubsection.\arabic{subsubsection}}

% --- Codigo fuente -----------------------------------------------------------
\usepackage{listings}
\lstdefinestyle{mystyle}{
    backgroundcolor=\color{backcolour},
    commentstyle=\color{codegreen},
    keywordstyle=\color{magenta},
    numberstyle=\tiny\color{codegray},
    stringstyle=\color{codepurple},
    basicstyle=\ttfamily\footnotesize,
    breakatwhitespace=false,
    breaklines=true,
    captionpos=b,
    keepspaces=true,
    numbers=left,
    numbersep=5pt,
    showspaces=false,
    showstringspaces=false,
    showtabs=false,
    tabsize=2
}
\lstset{style=mystyle}

% --- Diagramas TikZ ----------------------------------------------------------
\usepackage{tikz}
\usetikzlibrary{arrows.meta, shapes.geometric, positioning, fit, backgrounds, babel}

% --- Matematicas -------------------------------------------------------------
\usepackage{amsmath}
\usepackage{amssymb}
\usepackage{enumerate}

% --- Nombres localizados -----------------------------------------------------
\renewcommand{\figurename}{Fig.}
\renewcommand{\tablename}{Tabla}

\section{Ejercicio 1: Ajuste de controladores PID para sistemas SISO (S-IMC)}

Se dise\~nan los controladores PID para los dos procesos planteados en la Tarea no. 3,
utilizando la misma nomenclatura del ejemplo del m\'etodo S-IMC.

El modelo deseado se define como:
\[
G_f(s)=\frac{e^{-\theta s}}{\tau_c s+1},\qquad T_s \approx 4\tau_c+\theta
\]

Para el caso
\(G_p(s)=\dfrac{K_c e^{-\theta s}}{(\tau_1 s+1)(\tau_2 s+1)}\), se usa:
\[
K_p'=\frac{\tau_1}{K_c(\tau_c+\theta)},\qquad
T_i'=\min\left[\tau_1,\,4(\tau_c+\theta)\right],\qquad
T_d'=\tau_2
\]

Para el caso
\(G_p(s)=\dfrac{K_c e^{-\theta s}}{(\tau_1 s+1)}\), se usa:
\[
K_p'=\frac{1}{K_c(\tau_c+\theta)},\qquad
T_i'=\min\left[\tau_1,\,4(\tau_c+\theta)\right],\qquad
T_d'=0
\]

La conversi\'on de la estructura serie a formato ISA es:
\[
K_p = K_p'\left(1+\frac{T_d'}{T_i'}\right),\qquad
T_i = T_i'\left(1+\frac{T_d'}{T_i'}\right),\qquad
T_d = \frac{T_d'}{1+T_d'/T_i'}
\]

\subsection{Proceso 1}

\[
G_p(s)=\frac{0.4\,e^{-s}}{(5s+1)(2s+1)},\qquad T_s=19\,s
\]

Datos:
\[
K_c=0.4,\quad \theta=1,\quad \tau_1=5,\quad \tau_2=2
\]

  extbf{Paso 1: C\'alculo de \(\tau_c\)}
\[
19=4\tau_c+1 \Rightarrow \tau_c=4.5
\]

  extbf{Paso 2: Par\'ametros en estructura serie}
\[
K_p'=\frac{5}{0.4(4.5+1)}=2.2727
\]
\[
T_i'=\min\left[5,\,4(4.5+1)\right]=\min[5,22]=5
\]
\[
T_d'=2
\]

Resultado (serie):
\[
\boxed{K_p'=2.2727,\quad T_i'=5,\quad T_d'=2}
\]

\textbf{Paso 3: Conversi\'on a formato ISA}
\[
K_p=2.2727\left(1+\frac{2}{5}\right)=3.1818
\]
\[
T_i=5\left(1+\frac{2}{5}\right)=7
\]
\[
T_d=\frac{2}{1+2/5}=1.4286
\]

Resultado (ISA):
\[
\boxed{K_p=3.1818,\quad T_i=7,\quad T_d=1.4286}
\]

\subsection{Proceso 2}

\[
G_p(s)=\frac{1.3\,e^{-1.3s}}{(3.2s+1)},\qquad T_s=11\,s
\]

Datos:
\[
K_c=1.3,\quad \theta=1.3,\quad \tau_1=3.2
\]

\textbf{Paso 1: C\'alculo de \(\tau_c\)}
\[
11=4\tau_c+1.3 \Rightarrow \tau_c=2.425
\]

\textbf{Paso 2: Par\'ametros en estructura serie (PI)}
\[
K_p'=\frac{1}{1.3(2.425+1.3)}=\frac{1}{1.3(3.725)}=\frac{1}{4.8425}=0.2065
\]
\[
T_i'=\min\left[3.2,\,4(2.425+1.3)\right]=\min[3.2,14.9]=3.2
\]
\[
T_d'=0
\]

Resultado (serie):
\[
\boxed{K_p'=0.2065,\quad T_i'=3.2,\quad T_d'=0}
\]

\textbf{Paso 3: Conversi\'on a formato ISA}

Como \(T_d'=0\), los par\'ametros se mantienen:
\[
\boxed{K_p=0.2065,\quad T_i=3.2,\quad T_d=0}
\]

\section{Ejercicio 2: Controlador MIMO para sistema multitanques}

En este ejercicio se aplica la metodolog\'ia de dise\~no MIMO PID solicitada en la Tarea no. 3.
Se presenta el procedimiento completo y los resultados que deben calcularse a partir del
modelo del sistema multitanques de laboratorio.

\subsection{1) Modelo linealizado del sistema multitanques}

Se parte del modelo no lineal alrededor de un punto de operaci\'on
\((\mathbf{x}_0,\mathbf{u}_0,\mathbf{y}_0)\) y se linealiza mediante Jacobianos:
\[
\Delta\dot{\mathbf{x}}=A\,\Delta\mathbf{x}+B\,\Delta\mathbf{u},\qquad
\Delta\mathbf{y}=C\,\Delta\mathbf{x}+D\,\Delta\mathbf{u}
\]

La matriz de transferencia linealizada queda:
\[
G(s)=C(sI-A)^{-1}B+D
\]

Para un sistema 2\(\times\)2:
\[
G(s)=\begin{bmatrix}
G_{11}(s) & G_{12}(s)\\
G_{21}(s) & G_{22}(s)
\end{bmatrix}
\]

\subsection{2) Matriz de ganancias relativas (RGA)}

La RGA se calcula en estado estacionario con la ganancia est\'atica
\(K=G(0)\):
\[
\Lambda=K\circ\left(K^{-1}\right)^T
\]
donde \(\circ\) denota producto elemento a elemento.

Criterio de selecci\'on:
\begin{itemize}
  \item Preferir elementos de \(\Lambda\) cercanos a 1.
  \item Evitar pares con \(\lambda_{ij}<0\) o magnitudes muy grandes.
\end{itemize}

\subsection{3) Selecci\'on de lazos y criterio de emparejamiento}

Con base en \(\Lambda\), se define el emparejamiento entrada-salida que minimiza
la interacci\'on entre lazos. Para un caso t\'ipico 2\(\times\)2, el emparejamiento puede ser:
\[
y_1 \leftrightarrow u_1,\qquad y_2 \leftrightarrow u_2
\]
o, si la RGA lo indica, emparejamiento cruzado.

\subsection{4) Sinton\'ia IMC de cada lazo desacoplado}

Cada lazo principal se aproxima por FOPDT:
\[
G_{ii}(s)\approx \frac{K_i e^{-\theta_i s}}{\tau_i s+1}
\]

Usando IMC (forma PI):
\[
K_{p,i}=\frac{\tau_i}{K_i(\lambda_i+\theta_i)},\qquad
T_{i,i}=\tau_i
\]
donde \(\lambda_i\) es el par\'ametro de robustez deseado.

Si se requiere PID, se extiende con la forma utilizada en clase para SOPDT/FOPDT seg\'un el lazo.

\subsection{5) Simulaci\'on en Simulink con modelo no lineal}

Se implementa en Simulink:
\begin{itemize}
  \item Modelo no lineal de tanques.
  \item Dos controladores (uno por lazo) con el emparejamiento seleccionado.
  \item Pruebas de seguimiento de consigna y rechazo de perturbaciones.
\end{itemize}

Se reportan indicadores por lazo: tiempo de establecimiento, sobreimpulso,
error estacionario y esfuerzo de control.

\subsection{6) Estimaci\'on de flujos y arquitectura en cascada}

Se estima el flujo de enlace y/o salida a partir de mediciones de nivel y
caracter\'isticas hidr\'aulicas del sistema. Luego se plantea cascada:
\begin{itemize}
  \item Lazo interno r\'apido: control de flujo.
  \item Lazo externo lento: control de nivel.
\end{itemize}

Condici\'on recomendada de dise\~no en cascada:
\[
\omega_{c,\text{interno}} \approx 3\text{ a }5\;\omega_{c,\text{externo}}
\]

\subsection{7) Comparaci\'on de desempe\~no y reporte final}

Se compara el desempe\~no entre:
\begin{itemize}
  \item Control MIMO PID simple (sin cascada).
  \item Control con arquitectura en cascada.
\end{itemize}

La comparaci\'on se realiza para cambios de consigna y perturbaciones,
reportando m\'etricas de desempe\~no y an\'alisis de interacci\'on entre lazos.
\end{document}
